\slajdid{04-s034}
\begin{frame}[label=04-s034]{Komplement osnove}
% element: 04-s034:e01 — Definicija komplementa osnove i veza sa maksimumom+1 za r i n
% element: 04-s034:e02 — Binarni drugi komplement invertovanjem bita+1 i formula
% element: 04-s034:e03 — Decimalni primer: 1000 kodova, 0–499, −123=877 i podela vodećih cifara
% element: 04-s034:e04 — Izvorni naziv decimalnog primera: komplement 9-tke
Za osnovu $r$ i dužinu reči $n$ (broj cifara):
\[\alert{-V\equiv r^n-V=(r^n-1)-V+1}\]
U binarnom sistemu: \alert{komplement dvojke / drugi komplement}.\par
Dobija se invertovanjem bita i dodatkom $1$; „prirodan“ je za aritmetičke operacije.
{\fontsize{10}{12}\selectfont
\[111\ldots11-b_{n-1}b_{n-2}b_{n-3}\ldots b_1b_0+1=\overline{b_{n-1}}\,\overline{b_{n-2}}\,\overline{b_{n-3}}\ldots\overline{b_1}\,\overline{b_0}+1\]}
\NaslovBloka{Decimalni primer: tri cifre, 1000 kodova}
Pozitivne vrednosti: $0$–$499$. Prikaz $-123$ u komplementu 10-tke:
\[\alert{-123\equiv1000-123=877}\]
Vodeće cifre $0,1,2,3,4$: pozitivni; $5,6,7,8,9$: negativni brojevi.
\end{frame}
